\ifdefined\gkver\else\def\gkver{student}\fi
\documentclass[\gkver]{gaokaozhenti}
\begin{document}

\chapter{2004年江苏}

\begin{enumerate}
\item
%% number: 1
%% paperName: 2004年江苏高考第 1 题 4 分
%% typeId: 2
%% type: 多选题
%% chapter: 光学
%% point: 光子说
%% method: 无
%% score: 4
%% degree: 650
%% duplicateId: 0
%% body:
下列说法正确的是\xzanswer{AB}

\fourchoices[answer=AB]
{光波是一种概率波}
{光波是一种电磁波}
{单色光从光密介质进入光疏介质时，光子的能量改变}
{单色光从光密介质进入光疏介质时，光的波长不变}

%% answer: AB
%% memo:
\memoanswer{
	光波既是概率波，也是电磁波，故 A、B 正确；光子的能量由频率决定，在不同介质中保持不变，由 $\lambda=\frac{v}{f}$ 知，$v$ 由介质决定，故 $\lambda$ 也与介质有关，故 C、D 错误。
}

\item
%% number: 2
%% paperName: 2004年江苏高考第 2 题 4 分
%% typeId: 1
%% type: 单选题
%% chapter: 气体性质
%% point: 热力学第一定律
%% method: 无
%% score: 4
%% degree: 650
%% duplicateId: 0
%% body:
下列说法正确的是\xzanswer{D}

\fourchoices[answer=D]
{物体放出热量，温度一定降低}
{物体内能增加，温度一定升高}
{热量能自发地从低温物体传给高温物体}
{热量能自发地从高温物体传给低温物体}

%% answer: D
%% memo:
\memoanswer{
	由热力学第一定律 $\Delta U=W+Q$ 可知，只知 $Q$ 的变化，不能判断 $\Delta U$，故 A 错误；内能包括所有分子的动能和势能，温度是分子平均动能的标志，内能增加，温度不一定升高，故 B 错误；由热力学第二定律可知 C 错误、D 正确。
}

\item
%% number: 3
%% paperName: 2004年江苏高考第 3 题 4 分
%% typeId: 1
%% type: 单选题
%% chapter: 原子和原子核
%% point: 天然放射性
%% method: 无
%% score: 4
%% degree: 850
%% duplicateId: 0
%% body:
下列说法正确的是\xzanswer{C}

\fourchoices[answer=C]
{$\alpha$ 射线与 $\gamma$ 射线都是电磁波}
{$\beta$ 射线为原子的核外电子电离后形成的电子流}
{用加温、加压或改变其化学状态的方法都不能改变原子核衰变的半衰期}
{原子核经过衰变生成新核，则新核的质量总等于原核的质量}

%% answer: C
%% memo:
\memoanswer{
	$\alpha$、$\beta$、$\gamma$ 三种射线都是从原子核内发射出来的，$\alpha$ 射线是氦核流，$\beta$ 射线是高速电子流，$\gamma$ 射线是光子流（高频电磁波），故 A、B 错误；半衰期只由原子核的种类决定，故 C 正确；衰变过程有质量亏损，故 D 错误。
}

\item
%% number: 4
%% paperName: 2004年江苏高考第 4 题 4 分
%% typeId: 2
%% type: 多选题
%% chapter: 曲线运动
%% point: 人造地球卫星
%% method: 无
%% score: 4
%% degree: 650
%% duplicateId: 0
%% body:
若人造卫星绕地球作匀速圆周运动，则下列说法正确的是\xzanswer{BD}

\fourchoices[answer=BD]
{卫星的轨道半径越大，它的运行速度越大}
{卫星的轨道半径越大，它的运行速度越小}
{卫星的质量一定时，轨道半径越大，它需要的向心力越大}
{卫星的质量一定时，轨道半径越大，它需要的向心力越小}

%% answer: BD
%% memo:
\memoanswer{
	由万有引力提供向心力得 $G\frac{Mm}{r^{2}}=m\frac{v^{2}}{r}$，则 $v=\sqrt{\frac{GM}{r}}$，故 A 错误、B 正确；由 $F_{\text{向}}=G\frac{Mm}{r^{2}}$ 可知 C 错误、D 正确。
}

\item
%% number: 5
%% paperName: 2004年江苏高考第 5 题 4 分
%% typeId: 2
%% type: 多选题
%% chapter: 气体性质
%% point: 玻意耳定律
%% method: 无
%% score: 4
%% degree: 650
%% duplicateId: 0
%% body:
甲、乙两个相同的密闭容器中分别装有等质量的同种气体，已知甲、乙容器中气体的压强分别为 $p_{\text{甲}}$、$p_{\text{乙}}$，且 $p_{\text{甲}}<p_{\text{乙}}$，则\xzanswer{BC}

\fourchoices[answer=BC]
{甲容器中气体的温度高于乙容器中气体的温度}
{甲容器中气体的温度低于乙容器中气体的温度}
{甲容器中气体分子的平均动能小于乙容器中气体分子的平均动能}
{甲容器中气体分子的平均动能大于乙容器中气体分子的平均动能}

%% answer: BC
%% memo:
\memoanswer{
	由 $\frac{pV}{T}=$ 恒量，由 $p_{\text{甲}}<p_{\text{乙}}$，$V_{\text{甲}}=V_{\text{乙}}$ 可知 $T_{\text{甲}}<T_{\text{乙}}$，故 A 错误、B 正确；由于温度是分子平均动能的标志，故 C 正确、D 错误。
}

\item
%% number: 6
%% paperName: 2004年江苏高考第 6 题 4 分
%% typeId: 1
%% type: 单选题
%% chapter: 电磁感应
%% point: 线圈过有界磁场
%% method: 无
%% score: 4
%% degree: 750
%% duplicateId: 0
%% body:
如图所示，一个有界匀强磁场区域，磁场方向垂直纸面向外。一个矩形闭合导线框 abcd，沿纸面由位置 1（左）匀速运动到位置 2（右）。则\xzanswer{D}

\onepicture[w=5.5cm]{\includesvg{figs/06}}

\fourchoices[answer=D]
{导线框进入磁场时，感应电流方向为 $a\to b\to c\to d\to a$}
{导线框离开磁场时，感应电流方向为 $a\to d\to c\to b\to a$}
{导线框离开磁场时，受到的安培力方向水平向右}
{导线框进入磁场时，受到的安培力方向水平向左}

%% answer: D
%% memo:
\memoanswer{
	由楞次定律可知，进入磁场时感应电流的方向为顺时针方向，出磁场时感应电流的方向为逆时针方向，故 A、B 错误；由左手定则知，安培力方向均向左，故 C 错误、D 正确。
}

\item
%% number: 7
%% paperName: 2004年江苏高考第 7 题 4 分
%% typeId: 1
%% type: 单选题
%% chapter: 原子和原子核
%% point: 质能方程
%% method: 无
%% score: 4
%% degree: 650
%% duplicateId: 0
%% body:
雷蒙德•戴维斯因研究来自太阳的电子中微子（$\nu_{e}$）而获得了 2002 年度诺贝尔物理学奖。他探测中微子所用的探测器的主体是一个贮满 $615\Ut$ 四氯乙烯（\ce{C2Cl4}）溶液的巨桶，电子中微子可以将一个氯核转变为一个氩核，其核反应方程式为

$\ce{\nu_{e} + ^{37}_{17}Cl -> ^{37}_{18}Ar + ^{0}_{-1}e}$，

已知 \ce{^{37}_{17}Cl} 核的质量为 $36.95658\Uu$，\ce{^{37}_{18}Ar} 核的质量为 $36.95691\Uu$，\ce{^{0}_{-1}e} 的质量为 $0.00055\Uu$，$1\Uu$ 质量对应的能量为 $931.5\UMeV$。根据以上数据，可以判断参与上述反应的电子中微子的最小能量为\xzanswer{A}

\fourchoices[answer=A]
{$0.82\UMeV$}
{$0.31\UMeV$}
{$1.33\UMeV$}
{$0.51\UMeV$}

%% answer: A
%% memo:
\memoanswer{
	由题意有 $m_{\nu_{e}}+m_{\text{Cl}}\geq m_{\text{Ar}}+m_{\text{e}}$，得 $m_{\nu_{e}}\geq 0.00088\Uu$，则电子中微子的最小能量为 $0.00088\times 931.5\UMeV=0.82\UMeV$，故 A 正确。
}

\item
%% number: 8
%% paperName: 2004年江苏高考第 8 题 5 分
%% typeId: 2
%% type: 多选题
%% chapter: 机械波
%% point: 波与振动图象
%% method: 无
%% score: 5
%% degree: 450
%% duplicateId: 0
%% body:
图 \ref{2004江苏08a} 中，波源 S 从平衡位置 $y=0$ 开始振动，运动方向竖直向上（$y$ 轴的正方向），振动周期 $T=0.01\Us$，产生的简谐波向左、右两个方向传播，波速均为 $v=80\Ums$。经过一段时间后，P、Q 两点开始振动，已知距离 $SP=1.2\Um$、$SQ=2.6\Um$。若以 Q 点开始振动的时刻作为计时的零点，则在下列振动图象中，能正确描述 P、Q 两点振动情况的是\xzanswer{AD}

\onepicture[label=2004江苏08a, num=1, numstyle=arabic, showlabel=true, w=5.0cm]{\includesvg{figs/08a}}

\fourpicture[numstyle=chinese, labelA=2004江苏08b, labelB=2004江苏08c, labelC=2004江苏08d, labelD=2004江苏08e, w=3.4cm]
{figs/08b.png}
{figs/08c.png}
{figs/08d.png}
{figs/08e.png}

\fourchoices[answer=AD]
{甲为 Q 点振动图象}
{乙为 Q 点振动图象}
{丙为 P 点振动图象}
{丁为 P 点振动图象}

%% answer: AD
%% memo:
\memoanswer{
	由题意知波长 $\lambda=vT=0.8\Um$，由 $SQ-SP=1.4\Um=\frac{7}{4}\lambda$ 可知，P 点比 Q 点提前振动了 $\frac{7}{4}T$；若以 Q 点开始振动计时，则甲为 Q 点的振动图象，丁为 P 点的振动图象，故 A、D 正确。
}

\item
%% number: 9
%% paperName: 2004年江苏高考第 9 题 4 分
%% typeId: 1
%% type: 单选题
%% chapter: 光学
%% point: 全反射
%% method: 无
%% score: 4
%% degree: 650
%% duplicateId: 0
%% body:
如图所示，只含黄光和紫光的复色光束 PO，沿半径方向射入空气中的玻璃半圆柱后，被分成两光束 OA 和 OB 沿如图所示方向射出。则\xzanswer{C}

\onepicture[w=4.5cm]{figs/09.png}

\fourchoices[answer=C]
{OA 为黄光，OB 为紫光}
{OA 为紫光，OB 为黄光}
{OA 为黄光，OB 为复色光}
{OA 为紫光，OB 为复色光}

%% answer: C
%% memo:
\memoanswer{
	因 $n_{\text{紫}}>n_{\text{黄}}$，又 $n=\frac{1}{\sin C}$，故临界角 $C_{\text{紫}}<C_{\text{黄}}$，由图可知，光线 PO 的入射角恰好等于 $C_{\text{黄}}$，此时部分黄光沿 OA 射出，部分黄光和全部紫光沿 OB 射出，故 C 正确。
}

\item
%% number: 10
%% paperName: 2004年江苏高考第 10 题 4 分
%% typeId: 2
%% type: 多选题
%% chapter: 原子和原子核
%% point: 玻尔模型
%% method: 无
%% score: 4
%% degree: 550
%% duplicateId: 0
%% body:
若原子的某内层电子被电离形成空位，其它层的电子跃迁到该空位上时，会将多余的能量以电磁辐射的形式释放出来，此电磁辐射就是原子的特征 X 射线。内层空位的产生有多种机制，其中的一种称为内转换，即原子中处于激发态的核跃迁回基态时，将跃迁时释放的能量交给某一内层电子，使此内层电子电离而形成空位（被电离的电子称为内转换电子）。\ce{^{214}Po} 的原子核从某一激发态回到基态时，可将能量 $E_{0}=1.416\UMeV$ 交给内层电子（如 K、L、M 层电子，K、L、M 标记原子中最靠近核的三个电子层）使其电离。实验测得从 \ce{^{214}Po} 原子的 K、L、M 层电离出的电子的动能分别为 $E_{\text{k}}=1.323\UMeV$、$E_{\text{L}}=1.399\UMeV$、$E_{\text{M}}=1.412\UMeV$。则可能发射的特征 X 射线的能量为\xzanswer{AC}

\fourchoices[answer=AC]
{$0.013\UMeV$}
{$0.017\UMeV$}
{$0.076\UMeV$}
{$0.093\UMeV$}

%% answer: AC
%% memo:
\memoanswer{
	设能级为 $E_{n}$，电离后动能为 $E_{\text{动}}$，则 $E_{\text{动}}=E_{0}+E_{n}$，即 $E_{n}=E_{\text{动}}-E_{0}$，得 K、L、M 三个能级值 $E_{\text{K}}=-0.093\UMeV$、$E_{\text{L}}=-0.017\UMeV$、$E_{\text{M}}=-0.004\UMeV$，跃迁情形如图所示，共可发射三种特征 X 射线，能量分别为 $E_{1}=E_{\text{M}}-E_{\text{K}}=0.089\UMeV$、$E_{2}=E_{\text{L}}-E_{\text{K}}=0.076\UMeV$、$E_{3}=E_{\text{M}}-E_{\text{L}}=0.013\UMeV$，故 A、C 正确。

	\onepicture[w=3.6cm]{figs/10_an.jpg}
}

\item
%% number: 11
%% paperName: 2004年江苏高考第 11 题 8 分
%% typeId: 4
%% type: 实验题
%% chapter: 直线运动
%% point: 游标卡尺和螺旋测微仪
%% method: 无
%% score: 8
%% degree: 850
%% duplicateId: 0
%% body:
\begin{enumerate}
	\item
	某实验中需要测量一根钢丝的直径（约 $0.5\Umm$）。为了得到尽可能精确的测量数据，应从实验室提供的米尺、螺旋测微器和游标卡尺（游标尺上有 10 个等分刻度）中，选择 \tkanswer{螺旋测微器} 进行测量。

	\item
	用游标卡尺（游标尺上有 50 个等分刻度）测定某工件的宽度时，示数如图所示，此工件的宽度为 \tkanswer{23.22} $\Umm[a]$。
\end{enumerate}
\onepicture[align=right, w=6.5cm]{figs/11.png}

%% answer:
\jdanswer{
\begin{enumerate}
	\item
	螺旋测微器

	\item
	$23.22$
\end{enumerate}
}
%% memo:
\memoanswer{
\begin{enumerate}
	\item
	三种测量工具相比，螺旋测微器可以精确到 $0.001\Umm$，测量更为准确。

	\item
	游标卡尺读数为 $23\Umm+0.02\times 1\Umm=23.22\Umm$。
\end{enumerate}
}

\item
%% number: 12
%% paperName: 2004年江苏高考第 12 题 12 分
%% typeId: 4
%% type: 实验题
%% chapter: 电路
%% point: 伏安法测电阻
%% method: 无
%% score: 12
%% degree: 450
%% duplicateId: 0
%% body:
某同学对黑箱（见图 \ref{2004江苏12a}）中一个电学元件的伏安特性进行研究。通过正确测量，他发现该元件两端的电压 $U_{\text{ab}}$（$U_{\text{ab}}=U_{\text{a}}-U_{\text{b}}$）与流过它的电流 $I$ 之间的变化关系有如下规律：

①当 $-15\UV<U_{\text{ab}}<0\UV$ 时，$I$ 近似为零。

②当 $U_{\text{ab}}\geqslant 0$ 时，$U_{\text{ab}}$ 和 $I$ 的实验数据见下表：

\begin{center}
\begin{tabular}{|c|c|c|c|c|c|c|c|}
\hline
编号 & 1 & 2 & 3 & 4 & 5 & 6 & 7 \\
\hline
$U_{\text{ab}}/\UV$ & 0.000 & 0.245 & 0.480 & 0.650 & 0.705 & 0.725 & 0.745 \\
\hline
$I/\UmA$ & 0.00 & 0.15 & 0.25 & 0.60 & 1.70 & 4.25 & 7.50 \\
\hline
\end{tabular}
\end{center}

\onepicture[label=2004江苏12a, num=1, numstyle=arabic, showlabel=true, w=3.5cm]{figs/12a.png}

\begin{enumerate}
	\item
	在图 \ref{2004江苏12b} 中画出 $U_{\text{ab}}\geqslant 0$ 时该元件的伏安特性曲线。（可用铅笔作图）

	\onepicture[label=2004江苏12b, num=2, numstyle=arabic, showlabel=true, w=6.5cm]{figs/12b.png}

	\item
	根据上述实验事实，该元件具有的特性是 \tkanswer{单向导电性}。

	\item
	若将此黑箱接入图 \ref{2004江苏12c} 电路中，并在该电路的 cd 两端输入如图 \ref{2004江苏12d}（甲）所示的方波电压信号 $U_{\text{cd}}$，请在图 \ref{2004江苏12d}（乙）中定性画出负载电阻 $R_{L}$ 上的电压信号 $U_{\text{ef}}$ 的波形。

	\onepicture[label=2004江苏12c, num=3, numstyle=arabic, showlabel=true, w=5.5cm]{figs/12c.png}

	\onepicture[label=2004江苏12d, num=4, numstyle=arabic, showlabel=true, w=8cm]{figs/12d.png}
\end{enumerate}

%% answer:
\jdanswer{
\begin{enumerate}
	\item
	如图所示

	\item
	单向导电性

	\item
	如图所示
\end{enumerate}
}
%% memo:
\memoanswer{
\begin{enumerate}
	\item
	利用描点法，将各点坐标在图中描出，用平滑的曲线连接，如图所示。

	\onepicture[w=6.5cm]{figs/12_an1.png}

	\item
	由图可知，当 a 端电势小于 b 端时，元件中电流几乎为零，而反过来后，电流较大，故说明该元件具有单向导电性。

	\item
	若将该元件放入黑箱中，因为具有单向导电性，则电压正向时，有电流通过，此时输出端输出恒定电压，而电压反向时，无电流通过，输出端电压为零，故如图所示。

	\onepicture[w=8cm]{figs/12_an2.png}
\end{enumerate}
}

\item
%% number: 13
%% paperName: 2004年江苏高考第 13 题 14 分
%% typeId: 6
%% type: 计算题
%% chapter: 交流电
%% point: 变压器
%% method: 无
%% score: 14
%% degree: 850
%% duplicateId: 0
%% body:
如图所示，一个变压器（可视为理想变压器）的原线圈接在 $220\UV$ 的市电上，向额定电压为 $1.80\times 10^{4}\UV$ 的霓虹灯供电，使它正常发光。为了安全，需在原线圈回路中接入熔断器，使副线圈电路中电流超过 $12\UmA$ 时，熔丝就熔断。
\begin{enumerate}
	\item
	熔丝的熔断电流是多大？

	\item
	当副线圈电路中电流为 $10\UmA$ 时，变压器的输入功率是多大？
\end{enumerate}
\onepicture[align=right, w=5.5cm]{\includesvg{figs/13}}

%% answer:
\jdanswer{
\begin{enumerate}
	\item
	$I_{1}=0.98\UA$

	\item
	$P_{1}=180\UW$
\end{enumerate}
}
%% memo:
\memoanswer{
\begin{enumerate}
	\item
	设原、副线圈上的电压、电流分别为 $U_{1}$、$U_{2}$、$I_{1}$、$I_{2}$，由理想变压器的输入功率等于输出功率得
	$I_{1}U_{1}=I_{2}U_{2}$，
	当 $I_{2}=12\UmA$ 时，$I_{1}$ 即为熔断电流，代入数据得 $I_{1}=0.98\UA$。

	\item
	设副线圈中电流为 $I_{2}'=10\UmA$ 时，变压器的输入功率为 $P$，由理想变压器的输入功率等于输出功率得
	$P=I_{2}'U_{2}$，
	代入数据得 $P=180\UW$。
\end{enumerate}
}

\item
%% number: 14
%% paperName: 2004年江苏高考第 14 题 14 分
%% typeId: 6
%% type: 计算题
%% chapter: 电路
%% point: 电容
%% method: 无
%% score: 14
%% degree: 650
%% duplicateId: 0
%% body:
如图所示的电路中，电源电动势 $E=6.00\UV$，其内阻可忽略不计。电阻的阻值分别为 $R_{1}=2.4\UkO$、$R_{2}=4.8\UkO$，电容器的电容 $C=4.7\UuF$。闭合开关 S，待电流稳定后，用电压表测 $R_{1}$ 两端的电压，其稳定值为 $1.50\UV$。
\begin{enumerate}
	\item
	该电压表的内阻为多大？

	\item
	由于电压表的接入，电容器的带电量变化了多少？
\end{enumerate}
\onepicture[align=right, w=4.5cm]{figs/14.png}

%% answer:
\jdanswer{
\begin{enumerate}
	\item
	$R_{V}=4.8\UkO$

	\item
	$\Delta Q=2.35\times 10^{-6}\UC$
\end{enumerate}
}
%% memo:
\memoanswer{
\begin{enumerate}
	\item
	设电压表的内阻为 $R_{V}$，测得 $R_{1}$ 两端的电压为 $U_{1}$，$R_{1}$ 与 $R_{V}$ 并联后的总电阻为 $R$，则有
	$\frac{1}{R}=\frac{1}{R_{1}}+\frac{1}{R_{V}}$①，
	由串联电路的规律得
	$\frac{R}{R_{2}}=\frac{U_{1}}{E-U_{1}}$②，
	联立①②得
	$R_{V}=\frac{R_{1}R_{2}U_{1}}{R_{1}E-(R_{1}+R_{2})U_{1}}$，
	代入数据解得 $R_{V}=4.8\UkO$。

	\item
	电压表接入前，电容器上的电压 $U_{C}$ 等于电阻 $R_{2}$ 上的电压，$R_{1}$ 两端的电压为 $U_{R_{1}}$，则 $\frac{U_{C}}{U_{R_{1}}}=\frac{R_{2}}{R_{1}}$，又 $E=U_{C}+U_{R_{1}}$。
	接入电压表后，电容器上的电压为 $U_{C}'=E-U_{1}$。
	由于电压表的接入，电容器带电量增加了 $\Delta Q=C(U_{C}'-U_{C})$。
	由以上各式解得 $\Delta Q=C\left(\frac{R_{1}E}{R_{1}+R_{2}}-U_{1}\right)$，代入数据解得 $\Delta Q=2.35\times 10^{-6}\UC$。
\end{enumerate}
}

\item
%% number: 15
%% paperName: 2004年江苏高考第 15 题 15 分
%% typeId: 6
%% type: 计算题
%% chapter: 功和能
%% point: 整体机械能守恒
%% method: 无
%% score: 15
%% degree: 350
%% duplicateId: 0
%% body:
如图所示，半径为 $R$、圆心为 O 的大圆环固定在竖直平面内，两个轻质小圆环套在大圆环上。一根轻质长绳穿过两个小圆环，它的两端都系上质量为 $m$ 的重物，忽略小圆环的大小。
\begin{enumerate}
	\item
	将两个小圆环固定在大圆环竖直对称轴的两侧 $\theta=30^\circ$ 的位置上（如图）。在两个小圆环间绳子的中点 C 处，挂上一个质量 $M=\sqrt{2}m$ 的重物，使两个小圆环间的绳子水平，然后无初速释放重物 $M$。设绳子与大、小圆环间的摩擦均可忽略，求重物 $M$ 下降的最大距离。

	\item
	若不挂重物 $M$，小圆环可以在大圆环上自由移动，且绳子与大、小圆环间及大、小圆环之间的摩擦均可以忽略，问两个小圆环分别在哪些位置时，系统可处于平衡状态？
\end{enumerate}
\onepicture[align=right, h=6cm]{figs/15.png}

%% answer:
\jdanswer{
\begin{enumerate}
	\item
	$h=\sqrt{2}R$

	\item
	系统处于平衡状态时，两小环的可能位置为：a．两小环同时位于大圆环的底端；b．两小环同时位于大圆环的顶端；c．两小环一个位于大圆环的顶端，另一个位于大圆环的底端；d．除上述三种情况外，根据对称性可知，系统如能平衡，则两小圆环的位置一定关于大圆环竖直对称轴对称，设平衡时，两小圆环位于大圆环竖直对称轴两侧 $45^\circ$ 角的位置上。
\end{enumerate}
}
%% memo:
\memoanswer{
\begin{enumerate}
	\item
	重物向下先做加速运动，后做减速运动，当重物速度为零时，下降的距离最大。设下降的最大距离为 $h$，由机械能守恒定律得
	$Mgh=2mg\left[\sqrt{h^{2}+(R\sin\theta)^{2}}-R\sin\theta\right]$，
	又 $M=\sqrt{2}m$，解得 $h=\sqrt{2}R$（另解 $h=0$ 舍去）。

	\item
	系统处于平衡状态时，两小环的可能位置为：
	\begin{itemize}
		\item
		a．两小环同时位于大圆环的底端；
		\item
		b．两小环同时位于大圆环的顶端；
		\item
		c．两小环一个位于大圆环的顶端，另一个位于大圆环的底端；
		\item
		d．除上述三种情况外，由对称性可知，系统如能平衡，则两小圆环的位置一定关于大圆环竖直对称轴对称。设平衡时，两小圆环在大圆环竖直对称轴两侧 $\alpha$ 角的位置上（如图所示）。

		\onepicture[h=4.5cm]{figs/15_an.jpg}

		对于重物 $m$，受绳子拉力 $T$ 与重力 $mg$ 的作用，有 $T=mg$。
		对于小圆环，受到三个力的作用，水平绳子的拉力 $T$、竖直绳子的拉力 $T$、大圆环的支持力 $N$。两绳子的拉力沿大圆环切向的分力大小相等，方向相反，有 $T\sin\alpha=T\sin\alpha'$，解得 $\alpha=\alpha'$，又 $\alpha+\alpha'=90^\circ$，故 $\alpha=45^\circ$。
	\end{itemize}
\end{enumerate}
}

\item
%% number: 16
%% paperName: 2004年江苏高考第 16 题 15 分
%% typeId: 6
%% type: 计算题
%% chapter: 机械波
%% point: 多普勒效应
%% method: 无
%% score: 15
%% degree: 450
%% duplicateId: 0
%% body:
如图所示，声源 S 和观察者 A 都沿 $x$ 轴正方向运动，相对于地面的速率分别为 $v_{S}$ 和 $v_{A}$。空气中声音传播的速率为 $v_{p}$，设 $v_{S}<v_{p}$，$v_{A}<v_{p}$，空气相对于地面没有流动。
\begin{enumerate}
	\item
	若声源相继发出两个声信号，时间间隔为 $\Delta t$，请根据发出的这两个声信号从声源传播到观察者的过程，确定观察者接收到这两个声信号的时间间隔 $\Delta t'$；

	\item
	请利用（1）的结果，推导此情形下观察者接收到的声波频率与声源发出的声波频率间的关系式。
\end{enumerate}
\onepicture[align=right, w=8cm]{\includesvg{figs/16}}

%% answer:
\jdanswer{
\begin{enumerate}
	\item
	$\Delta t'=\frac{v_{p}-v_{S}}{v_{p}-v_{A}}\Delta t$

	\item
	$f'=\frac{v_{p}-v_{A}}{v_{p}-v_{S}}f$
\end{enumerate}
}
%% memo:
\memoanswer{
\begin{enumerate}
	\item
	设 $t_{1}$、$t_{2}$ 为声源 S 发出两个信号的时刻，$t_{1}'$、$t_{2}'$ 为观察者接收到两个信号的时刻，则第一个信号经过 $(t_{1}'-t_{1})$ 时间被观察者 A 接收到，第二个信号经过 $(t_{2}'-t_{2})$ 时间被观察者 A 接收到，且 $t_{2}-t_{1}=\Delta t$，$t_{2}'-t_{1}'=\Delta t'$。
	设声源发出第一个信号时，S、A 两点间的距离为 $L$，两个声信号从声源传播到观察者的过程中，它们运动的距离关系如图所示，可得
	$v_{p}(t_{1}'-t_{1})=L+v_{A}(t_{1}'-t_{1})$，
	$v_{p}(t_{2}'-t_{2})=L+v_{A}(t_{2}'-t_{1})-v_{S}\Delta t$，
	由以上各式，得 $\Delta t'=\frac{v_{p}-v_{S}}{v_{p}-v_{A}}\Delta t$。

	\onepicture[w=7cm]{figs/16_an1.jpg}

	\onepicture[w=7cm]{figs/16_an2.jpg}

	\item
	设声源发出声波的振动周期为 $T$，这样，由以上结论，观察者接收到的声波振动的周期 $T'=\frac{v_{p}-v_{S}}{v_{p}-v_{A}}T$，由此可得，观察者接收到的声波频率与声源发出声波频率间的关系为 $f'=\frac{v_{p}-v_{A}}{v_{p}-v_{S}}f$。
\end{enumerate}
}

\item
%% number: 17
%% paperName: 2004年江苏高考第 17 题 16 分
%% typeId: 6
%% type: 计算题
%% chapter: 电场
%% point: 带电粒子的偏转
%% method: 无
%% score: 16
%% degree: 450
%% duplicateId: 0
%% body:
汤姆生用来测定电子的比荷（电子的电荷量与质量之比）的实验装置如图所示，真空管内的阴极 K 发出的电子（不计初速、重力和电子间的相互作用）经加速电压加速后，穿过 $A'$ 中心的小孔沿中心轴 $O_{1}O$ 的方向进入到两块水平正对放置的平行极板 P 和 $P'$ 间的区域。当极板间不加偏转电压时，电子束打在荧光屏的中心 O 点处，形成了一个亮点；加上偏转电压 $U$ 后，亮点偏离到 $O'$ 点，$O'$ 与 O 点的竖直间距为 $d$，水平间距可忽略不计。此时，在 P 和 $P'$ 间的区域，再加上一个方向垂直于纸面向里的匀强磁场。调节磁场的强弱，当磁感应强度的大小为 $B$ 时，亮点重新回到 O 点。已知极板水平方向的长度为 $L_{1}$，极板间距为 $b$，极板右端到荧光屏的距离为 $L_{2}$（如图所示）。
\begin{enumerate}
	\item
	求打在荧光屏 O 点的电子速度的大小。

	\item
	推导出电子的比荷的表达式。
\end{enumerate}
\onepicture[align=right, w=8cm]{figs/17.png}

%% answer:
\jdanswer{
\begin{enumerate}
	\item
	$v=\frac{U}{Bb}$

	\item
	$\frac{e}{m}=\frac{Ud}{B^{2}bL_{1}\left(L_{2}+\frac{L_{1}}{2}\right)}$
\end{enumerate}
}
%% memo:
\memoanswer{
\begin{enumerate}
	\item
	当电子受到的电场力与洛伦兹力平衡时，电子做匀速直线运动，亮点重新回复到中心 O 点，设电子的速度为 $v$，则 $evB=eE$，解得 $v=\frac{E}{B}$，又 $U=Eb$，即 $v=\frac{U}{Bb}$。

	\item
	当极板间仅有偏转电场时，电子以速度 $v$ 进入后，竖直方向做匀加速运动，加速度大小为 $a=\frac{eU}{mb}$，电子在水平方向做匀速运动，在电场内的运动时间为 $t_{1}=\frac{L_{1}}{v}$，这样，电子在电场中，竖直向上偏转的距离为 $d_{1}=\frac{1}{2}at_{1}^{2}=\frac{eL_{1}^{2}U}{2mv^{2}b}$，离开电场时竖直向上的分速度为 $v_{\perp}=at_{1}=\frac{eL_{1}U}{mvb}$，电子离开电场后做匀速直线运动，经 $t_{2}$ 时间到达荧光屏，有 $t_{2}=\frac{L_{2}}{v}$，$t_{2}$ 时间内向上运动的距离为 $d_{2}=v_{\perp}t_{2}=\frac{eUL_{1}L_{2}}{mv^{2}b}$，这样，电子向上的总偏转距离为 $d=d_{1}+d_{2}=\frac{eU}{mv^{2}b}L_{1}\left(L_{2}+\frac{L_{1}}{2}\right)$，可解得 $\frac{e}{m}=\frac{Ud}{B^{2}bL_{1}\left(L_{2}+\frac{L_{1}}{2}\right)}$。
\end{enumerate}
}

\item
%% number: 18
%% paperName: 2004年江苏高考第 18 题 14 分
%% typeId: 6
%% type: 计算题
%% chapter: 运动和力
%% point: 动量守恒定律
%% method: 无
%% score: 14
%% degree: 350
%% duplicateId: 0
%% body:
一个质量为 $M$ 的雪橇静止在水平雪地上，一条质量为 $m$ 的爱斯基摩狗站在该雪橇上。狗向雪橇的正后方跳下，随后又追赶并向前跳上雪橇；其后狗又反复地跳下、追赶并跳上雪橇，狗与雪橇始终沿一条直线运动。若狗跳离雪橇时雪橇的速度为 $v$，则此时狗相对于地面的速度为 $v+u$（其中 $u$ 为狗相对于雪橇的速度，$v+u$ 为代数和。若以雪橇运动的方向为正方向，则 $v$ 为正值，$u$ 为负值）。设狗总以速度 $v$ 追赶和跳上雪橇，雪橇与雪地间的摩擦忽略不计。已知 $v$ 的大小为 $5\Ums$，$u$ 的大小为 $4\Ums$，$M=30\Ukg$，$m=10\Ukg$。
\begin{enumerate}
	\item
	求狗第一次跳上雪橇后两者的共同速度的大小。

	\item
	求雪橇最终速度的大小和狗最多能跳上雪橇的次数。
\end{enumerate}

（供使用但不一定用到的对数值：$\lg 2=0.301$，$\lg 3=0.477$）

%% answer:
\jdanswer{
\begin{enumerate}
	\item
	$V_{1}'=2\Ums$

	\item
	$V_{4}=5.625\Ums$，狗最多能跳上雪橇 3 次。
\end{enumerate}
}
%% memo:
\memoanswer{
\begin{enumerate}
	\item
	设雪橇运动的方向为正方向，狗第 1 次跳下雪橇后，雪橇的速度为 $V_{1}$，由动量守恒得 $MV_{1}+m(V_{1}+u)=0$，狗第 1 次跳上雪橇时，雪橇与狗的共同速度为 $V_{1}'$，满足 $MV_{1}+mv=(M+m)V_{1}'$，可解得 $V_{1}'=\frac{-Mmu+(M+m)mv}{(M+m)^{2}}$，将 $u=-4\Ums$、$v=5\Ums$、$M=30\Ukg$、$m=10\Ukg$ 代入解得 $V_{1}'=2\Ums$。

	\item
	解法一：设雪橇运动的方向为正方向，狗第 $(n-1)$ 次跳下雪橇后雪橇的速度为 $V_{n-1}$，则狗第 $(n-1)$ 次跳上雪橇后的速度 $V_{n-1}'$，满足 $MV_{n-1}+mv=(M+m)V_{n-1}'$。这样，狗 $n$ 次跳下雪橇后，雪橇的速度 $V_{n}$ 满足 $(M+m)V_{n-1}'=MV_{n}+m(V_{n}+u)$，解得
	$V_{n}=(v-u)\left[1-\left(\frac{M}{M+m}\right)^{n-1}\right]-\frac{mu}{M+m}\left(\frac{M}{M+m}\right)^{n-1}$，
	狗追不上雪橇的条件是 $V_{n}\geqslant v$，可化简为 $\left(\frac{M}{M+m}\right)^{n-1}\leqslant\frac{(M+m)u}{Mu-(M+m)v}$，最后可求得
	$n\geqslant 1+\frac{\lg\left(\frac{(M+m)u}{Mu-(M+m)v}\right)}{\lg\left(\frac{M}{M+m}\right)}$，
	代入数据解得 $n\geqslant 3.41$，故狗最多能跳上雪橇 3 次，雪橇最终的速度大小为 $V_{4}=5.625\Ums$。

	解法二：设雪橇运动的方向为正方向。狗第 $i$ 次跳下雪橇后，雪橇的速度为 $V_{i}$，狗的速度为 $V_{i}+u$；狗第 $i$ 次跳上雪橇后，雪橇和狗的共同速度为 $V_{i}'$，由动量守恒定律得
	第一次跳下雪橇：$MV_{1}+m(V_{1}+u)=0$，解得 $V_{1}=-\frac{mu}{M+m}=1\Ums$；
	第一次跳上雪橇：$MV_{1}+mv=(M+m)V_{1}'$；
	第二次跳下雪橇：$(M+m)V_{1}'=MV_{2}+m(V_{2}+u)$，解得 $V_{2}=\frac{(M+m)V_{1}'-mu}{M+m}=3\Ums$；
	第二次跳上雪橇：$MV_{2}+mv=(M+m)V_{2}'$，解得 $V_{2}'=\frac{MV_{2}+mv}{M+m}$；
	第三次跳下雪橇：$(M+m)V_{2}'=MV_{3}+m(V_{3}+u)$，解得 $V_{3}=\frac{(M+m)V_{2}'-mu}{M+m}=4.5\Ums$；
	第三次跳上雪橇：$MV_{3}+mv=(M+m)V_{3}'$，解得 $V_{3}'=\frac{MV_{3}+mv}{M+m}$；
	第四次跳下雪橇：$(M+m)V_{3}'=MV_{4}+m(V_{4}+u)$，解得 $V_{4}=\frac{(M+m)V_{3}'-mu}{M+m}=5.625\Ums$。
	此时雪橇的速度已大于狗追赶的速度，狗将不可能追上雪橇，因此，狗最多能跳上雪橇 3 次，雪橇最终的速度大小为 $5.625\Ums$。
\end{enumerate}
}

\end{enumerate}

\end{document}
